\section{Построение теней}

Построение теней сводится к удалению невидимых поверхностей для положения источника~\cite[с.~417]{rogers1989}. Свечение сияния "--- протяжённый источник, меняющийся во времени, поэтому метод должен применяться к нескольким источникам в каждом кадре.

\subsection{Метод теневого z-буфера}

Сцена изображается из положения источника, и глубины видимых точек сохраняются в теневом буфере $Z_{\text{т}}$. При построении кадра видимая точка $P$ переводится в систему координат источника $(x', y', z')$ и признак её освещённости равен~\cite[с.~422--423]{rogers1989}
\begin{equation}
	S(P) =
	\begin{cases}
		1, & z'_P \leqslant Z_{\text{т}}(x'_P, y'_P), \\
		0, & \text{иначе};
	\end{cases}
	\label{eq:shadow-test}
\end{equation}
вклад источника в формуле~\eqref{eq:blinn-phong} умножается на $S(P)$. Нужен один дополнительный проход на источник; конечное разрешение буфера даёт ступенчатость границ теней~\cite[с.~423]{rogers1989}.

\subsection{Трассировка теневых лучей}

Из видимой точки $P$ строится луч к источнику
\begin{equation}
	P(t) = P + t \vec{L}, \qquad t > 0;
	\label{eq:shadow-ray}
\end{equation}
если он пересекает объект, точка находится в тени~\cite[с.~424--425]{rogers1989}. Метод точен, но без ускоряющих структур данных требует порядка $p \cdot n$ проверок на источник, где $p$ "--- число пикселей кадра, $n$ "--- число граней сцены.

\subsection{Выбор метода построения теней}

Оценки в таблице~\ref{tab:shadows} получены для ландшафта из $n \approx 1{,}3 \cdot 10^5$ треугольников (карта высот $256 \times 256$), кадра $1920 \times 1080$ ($p \approx 2{,}1 \cdot 10^6$) и теневого буфера $2048 \times 2048$ ($p_s \approx 4{,}2 \cdot 10^6$) как $n + p_s + p$ и $p \cdot n$ действий соответственно.

\begin{table}[H]
	\caption{Сравнение методов построения теней}
	\label{tab:shadows}
	\begin{tabular}{|L{0.44}|C{0.27}|C{0.27}|}
		\hline
		Критерий & Теневой z-буфер & Теневые лучи \\
		\hline
		Действий на кадр для одного источника & $\approx 6{,}4 \cdot 10^{6}$ & $\approx 2{,}7 \cdot 10^{11}$ \\
		\hline
		Пересчёт при движении наблюдателя & не нужен & нужен \\
		\hline
		Искажения дискретизации & есть & нет \\
		\hline
	\end{tabular}
\end{table}

Теневой z-буфер требует примерно в $4 \cdot 10^4$ раз меньше действий на каждый источник, а его искажения уменьшаются увеличением размера буфера, поэтому выбран он.